% Nội dung mục 4.5 -- tạo từ kết quả trong anova_analysis/results.
% Không đặt \documentclass hay \begin{document} trong tệp thành phần.
\section{Kết quả ANOVA, kiểm tra giả định và hậu kiểm}
\label{sec:ketqua-anova}

Phần này áp dụng quy trình tại Mục~\ref{sec:phuongphap-anova} cho bốn câu hỏi nghiên
cứu. Kết quả chính dùng Welch ANOVA; hậu kiểm dùng Games--Howell và bốn p-value tổng thể
được hiệu chỉnh Holm. Các bảng mô tả vẫn giữ đơn vị gốc để hỗ trợ diễn giải thực tế.

\subsection{Thống kê mô tả và kiểm tra giả định}

Bảng~\ref{tab:anova:descriptive} cho thấy trung bình và trung vị khác nhau đáng kể ở
nhiều nhóm. Đặc biệt, trung bình \texttt{src\_bytes} của Probe đạt 385.679,84 byte nhưng
trung vị chỉ là 1 byte; \texttt{duration} của Probe có trung bình 2.074,86 giây nhưng
trung vị bằng 0. Đây là dấu hiệu đuôi phải dài và khối lượng lớn tại 0, không phải bằng
chứng để tùy tiện loại ngoại lệ.

\begin{table}[H]
\centering
\caption{Trung bình và trung vị theo nhóm trên thang gốc}
\label{tab:anova:descriptive}
\scriptsize
\setlength{\tabcolsep}{3pt}
\resizebox{\textwidth}{!}{%
\begin{tabular}{llrrrrr}
\toprule
Đặc trưng & Thống kê & Normal & DoS & Probe & R2L & U2R \\
\midrule
\texttt{count} & Trung bình & 22,52 & 178,09 & 77,05 & 1,30 & 5,81 \\
                & Trung vị   & 4 & 172 & 1 & 1 & 1 \\
\texttt{serror\_rate} & Trung bình & 0,013 & 0,748 & 0,047 & 0,012 & 0,038 \\
                & Trung vị   & 0 & 1 & 0 & 0 & 0 \\
\texttt{src\_bytes} & Trung bình & 13.133,28 & 1.176,32 & 385.679,84 & 307.727,30 & 906,23 \\
                & Trung vị   & 233 & 0 & 1 & 334 & 273 \\
\texttt{duration} & Trung bình & 168,59 & 0,006 & 2.074,86 & 633,42 & 80,94 \\
                & Trung vị   & 0 & 0 & 0 & 0 & 46 \\
\bottomrule
\end{tabular}}
\par\smallskip
\footnotesize Cỡ mẫu theo thứ tự Normal, DoS, Probe, R2L, U2R là 67.343, 45.927,
11.656, 995 và 52. \texttt{count}: kết nối; \texttt{src\_bytes}: byte;
\texttt{duration}: giây; \texttt{serror\_rate}: tỷ lệ 0--1.
\end{table}

Brown--Forsythe bác bỏ đồng nhất phương sai ở cả bốn đặc trưng
(Bảng~\ref{tab:anova:assumptions}); tỷ số cỡ mẫu lớn nhất/nhỏ nhất đạt 1.295,1. ECDF ở
Hình~\ref{fig:anova:ecdf} và Q--Q plot ở Hình~\ref{fig:anova:qq} đồng thời cho thấy dữ
liệu bị chặn, nhiều số 0 hoặc có đuôi dài. Các dấu hiệu này, thay vì chỉ một p-value
Shapiro--Wilk, là cơ sở chọn Welch và thang $\ln(1+x)$ cho hai biến lệch mạnh.

\begin{table}[H]
\centering
\caption{Kiểm tra giả định trên thang dùng cho kiểm định chính}
\label{tab:anova:assumptions}
\small
\setlength{\tabcolsep}{4pt}
\begin{tabularx}{\textwidth}{l l r l X}
\toprule
Đặc trưng & Thang phân tích & Brown--Forsythe & p-value & Chẩn đoán trọng tâm \\
\midrule
\texttt{count} & Gốc & 5.817,05 & $<10^{-300}$ & Biến đếm; $|\mathrm{skew}|$ tối đa 5,56; U2R nhỏ. \\
\texttt{serror\_rate} & Gốc & 5.279,29 & $<10^{-300}$ & Bị chặn 0--1, dồn tại hai biên; $|\mathrm{skew}|$ tối đa 8,85. \\
\texttt{src\_bytes} & $\ln(1+x)$ & 364,36 & $1{,}59\times10^{-312}$ & Thang gốc có cực trị lớn; sau log vẫn chưa chuẩn hoàn toàn. \\
\texttt{duration} & $\ln(1+x)$ & 1.318,15 & $<10^{-300}$ & Nhiều số 0; $|\mathrm{skew}|$ tối đa 38,03 do nhóm DoS. \\
\bottomrule
\end{tabularx}
\end{table}

\begin{figure}[H]
\centering
\includegraphics[width=0.98\linewidth]{../ANOVA/figures/anova_ecdf_four_features.pdf}
\caption{ECDF theo năm nhóm. \texttt{src\_bytes} và \texttt{duration} dùng
$\ln(1+x)$; lấy mẫu cố định tối đa 5.000 quan sát/nhóm chỉ để vẽ theo quy trình ở
Mục~\ref{sec:phuongphap-anova}. Mọi kiểm định sử dụng toàn bộ dữ liệu.}
\label{fig:anova:ecdf}
\end{figure}

\begin{figure}[H]
\centering
\includegraphics[width=0.98\linewidth]{../ANOVA/figures/qq_log_features.pdf}
\caption{Q--Q plot của $\log(1+\texttt{src\_bytes})$ và
$\log(1+\texttt{duration})$, dùng tối đa 1.000 quan sát/nhóm. Các điểm lệch đường tham
chiếu cho thấy phép log giảm nhưng không loại bỏ hoàn toàn sai lệch chuẩn.}
\label{fig:anova:qq}
\end{figure}

\subsection{Kết quả kiểm định tổng thể và kích thước ảnh hưởng}

Bảng~\ref{tab:anova:omnibus} bác bỏ $H_0$ cho cả bốn câu hỏi, kể cả sau hiệu chỉnh Holm.
Các p-value dưới giới hạn hiển thị được báo cáo là $p<10^{-300}$, không viết $p=0$.
Kết quả của \texttt{src\_bytes} và \texttt{duration} áp dụng cho trung bình trên thang
$\ln(1+x)$, không phải chênh lệch byte hoặc giây trung bình trên thang gốc.

\begin{table}[H]
\centering
\caption{Kết quả Welch ANOVA cho bốn câu hỏi nghiên cứu}
\label{tab:anova:omnibus}
\small
\setlength{\tabcolsep}{3.5pt}
\resizebox{\textwidth}{!}{%
\begin{tabular}{llllrrrr}
\toprule
RQ & Biến phân tích & $F$ & Bậc tự do $(df_1;df_2)$ & p thô & p Holm & $\eta^2_{\mathrm{mô\ tả}}$ & $\omega^2_{\mathrm{mô\ tả}}$ \\
\midrule
RQ1 & \texttt{count} & 35.907,45 & (4; 405,490) & $<10^{-300}$ & $<10^{-300}$ & 0,4049 & 0,4049 \\
RQ2 & \texttt{serror\_rate} & 31.976,20 & (4; 386,813) & $<10^{-300}$ & $<10^{-300}$ & 0,6202 & 0,6202 \\
RQ3 & $\ln(1+\texttt{src\_bytes})$ & 42.494,96 & (4; 387,016) & $<10^{-300}$ & $<10^{-300}$ & 0,5639 & 0,5639 \\
RQ4 & $\ln(1+\texttt{duration})$ & 1.652,44 & (4; 386,205) & $2{,}16\times10^{-241}$ & $2{,}16\times10^{-241}$ & 0,0412 & 0,0412 \\
\bottomrule
\end{tabular}}
\par\smallskip
\footnotesize $\eta^2$ và $\omega^2$ được tính từ tổng bình phương cổ điển trên thang
phân tích, nên là chỉ số mô tả, không phải effect size suy ra trực tiếp từ Welch $F$.
Không dùng các giá trị này để xếp hạng khả năng dự đoán của đặc trưng.
\end{table}

Effect size cho thấy p-value rất nhỏ không đồng nghĩa mức phân tách thực tiễn giống nhau:
$\eta^2$ mô tả của \texttt{serror\_rate} là 0,6202, trong khi \texttt{duration} chỉ là
0,0412 trên thang phân tích.

\subsection{Hậu kiểm Games--Howell và trả lời câu hỏi nghiên cứu}

Games--Howell kiểm soát so sánh nhiều cặp trong từng RQ và không giả định phương sai hay
cỡ mẫu bằng nhau. Bảng~\ref{tab:anova:posthoc-summary} tóm tắt toàn bộ 10 cặp/RQ;
Hình~\ref{fig:anova:posthoc} trình bày p-value hiệu chỉnh bằng cả số và ký hiệu, nên
không cần suy luận chỉ dựa trên màu.

\begin{table}[H]
\centering
\caption{Tóm tắt hậu kiểm Games--Howell trong từng câu hỏi}
\label{tab:anova:posthoc-summary}
\small
\setlength{\tabcolsep}{4pt}
\begin{tabularx}{\textwidth}{l c X}
\toprule
RQ & Cặp có ý nghĩa & Các cặp chưa đủ bằng chứng khác biệt ở mức 5\% \\
\midrule
RQ1 -- \texttt{count} & 9/10 & R2L--U2R: chênh lệch $-4{,}51$, KTC 95\% $[-13{,}71;4{,}69]$, $p_{adj}=0{,}639$. \\
RQ2 -- \texttt{serror\_rate} & 6/10 & Normal--R2L; Normal--U2R; Probe--U2R; R2L--U2R. \\
RQ3 -- $\ln(1+\texttt{src\_bytes})$ & 9/10 & Normal--U2R: chênh lệch 0,455, KTC 95\% $[-0{,}733;1{,}643]$, $p_{adj}=0{,}815$. \\
RQ4 -- $\ln(1+\texttt{duration})$ & 10/10 & Không có; cả 10 khoảng tin cậy đồng thời không chứa 0. \\
\bottomrule
\end{tabularx}
\end{table}

\begin{figure}[H]
\centering
\includegraphics[width=0.92\linewidth]{../ANOVA/figures/games_howell_heatmaps.pdf}
\caption{P-value Games--Howell đã hiệu chỉnh trong 10 so sánh cặp của từng RQ.
Dấu * biểu thị $p_{adj}<0{,}05$; đường viền trắng phân tách từng ô. Thanh màu bên phải
biểu diễn p-value đã hiệu chỉnh trong miền 0--0,05; giá trị từ 0,05 trở lên được hiển thị
tại cận trên. Ô không có dấu * không được diễn giải là hai trung bình bằng nhau, mà chỉ
là chưa đủ bằng chứng khác biệt với dữ liệu hiện có.}
\label{fig:anova:posthoc}
\end{figure}

\paragraph{RQ1 -- số lượng kết nối gần đây.}
Trung bình \texttt{count} khác nhau giữa ít nhất hai nhóm. DoS cao nhất (178,09), tiếp
đến Probe (77,05) và Normal (22,52); R2L và U2R lần lượt là 1,30 và 5,81. Games--Howell
cho thấy 9/10 cặp khác biệt, ngoại trừ R2L--U2R. Kết quả ủng hộ nhận định hành vi tạo kết
nối trong cửa sổ hai giây khác nhau rõ giữa DoS, Probe và Normal, nhưng không chứng minh
mọi dạng tấn công đều có \texttt{count} cao.

\paragraph{RQ2 -- tỷ lệ lỗi SYN.}
DoS có trung bình \texttt{serror\_rate}=0,748 và trung vị 1, vượt xa Normal
(trung bình 0,013; trung vị 0). Chênh lệch DoS--Normal là 0,735 với KTC đồng thời 95\%
$[0{,}729;0{,}741]$, $p_{adj}<10^{-300}$. Tuy vậy, chỉ 6/10 cặp có ý nghĩa và phân bố
dồn mạnh ở 0 hoặc 1. ANOVA vì thế cung cấp xấp xỉ hữu ích về trung bình tỷ lệ trong mẫu,
không phải một mô hình đầy đủ cho biến tỷ lệ bị dồn biên.

\paragraph{RQ3 -- lượng byte gửi từ nguồn.}
Trên thang gốc, Probe và R2L có trung bình rất cao (385.679,84 và 307.727,30 byte) nhưng
trung vị chỉ là 1 và 334 byte, cho thấy trung bình bị chi phối bởi đuôi phải. Welch ANOVA
trên $\ln(1+\texttt{src\_bytes})$ bác bỏ $H_0$; 9/10 cặp có ý nghĩa, riêng Normal--U2R
chưa đủ bằng chứng. Phát hiện này trả lời câu hỏi về trung bình sau biến đổi log, không
được chuyển thành phát biểu rằng một nhóm chênh một số byte trung bình tương ứng.

\paragraph{RQ4 -- thời gian kết nối.}
Welch ANOVA trên $\ln(1+\texttt{duration})$ bác bỏ $H_0$ và cả 10 cặp Games--Howell có
ý nghĩa. Tuy nhiên, trung vị bằng 0 ở bốn nhóm và tỷ lệ bằng 0 thay đổi từ 28,85\% đến
99,91\% (Hình~\ref{fig:anova:duration-zero}), nên kết quả chỉ được xem là bằng chứng
tổng hợp trên thang log, không phải một mô hình hoàn chỉnh cho dữ liệu zero-inflated.

\begin{figure}[H]
\centering
\includegraphics[width=0.88\linewidth]{../ANOVA/figures/duration_zero_wilson.pdf}
\caption{Tỷ lệ \texttt{duration}=0 và khoảng tin cậy Wilson 95\% theo nhóm. Đây là
phần nhị phân của phân tích hai phần, không phải khoảng tin cậy của thời lượng trung bình.}
\label{fig:anova:duration-zero}
\end{figure}

Trong 10.018 kết nối có \texttt{duration}$>0$, phân tích bổ sung bằng Welch trên
$\ln(1+\texttt{duration})$ vẫn có ý nghĩa,
$F(4;158{,}324)=315{,}82$, $p=2{,}48\times10^{-74}$. Đây là kết quả có điều kiện và
không thay thế kết quả trên toàn bộ mẫu.

\noindent\textbf{Ý nghĩa khoa học.}
ANOVA cho thấy bốn khía cạnh hành vi mạng có phân bố trung bình khác nhau theo nhóm,
nhưng mức chồng lấp, thang đo và độ ổn định của U2R khác nhau đáng kể. Kết quả hỗ trợ
diễn giải những tín hiệu mà mô hình học máy có thể khai thác, song p-value nhỏ không đo
trực tiếp năng lực dự đoán và không chứng minh quan hệ nhân quả. Do NSL-KDD là bộ chuẩn
lịch sử và không có định danh phiên đủ để kiểm chứng độc lập, suy luận chỉ áp dụng cho
mẫu đang phân tích; Welch xử lý phương sai không bằng nhau nhưng không khắc phục được
phụ thuộc tiềm ẩn giữa các kết nối.

\FloatBarrier

